% =========================
% 6_experiments.tex  --- protocol and tables written; result cells reserved (\PH / \PENDING)
% =========================
\section{Experiments}
\label{sec:experiments}

\PENDING{Corpus, protocol, and table structure are final; orange cells are reserved for results.}

\subsection{Methodology}
\label{sec:methodology}

We ask four questions. \textbf{(Q1)} Along the label-efficiency curve, where does HALO stand
against strong frozen encoders, and is there a regime it owns? \textbf{(Q2)} Does the tokenizer's
inductive bias do what Section~\ref{sec:no_grammar} claims---and, decisively, does anchoring let
one flat memory serve devices it was not built from? \textbf{(Q3)} Which Phase-A objectives carry
the transfer? \textbf{(Q4)} Does the system run on a phone?

\subsubsection{Data}
\label{sec:datasets}

We train on 12 public HAR datasets contributing 20 distinct sensor streams and evaluate on 7
never seen during training; Appendix~\ref{app:corpus} lists both with per-dataset rates,
placements, and window counts, and records the six datasets deliberately excluded and why. After
conversion to native-rate grids and removal of 95 windows exceeding physically possible sensor
rails, the corpus is $343{,}140$ non-overlapping six-second windows---$572$ hours---over $93$
labels, split subject-disjointly into $300{,}231$ training and $42{,}909$ validation windows, the
latter built by a greedy label-cover procedure so all $93$ labels are represented. Training rates
span $20$--$100$\,Hz over placements at the wrist, waist, pocket, lower back, ear, and head; the
held-out sets range from $1{,}260$ to $11{,}237$ windows over $6$ to $13$ labels.

\subsubsection{Baselines and the fairness contract}
\label{sec:baselines}

Cross-model comparison here is easy to get wrong, because the models do not accept the same
inputs. We state one contract and apply it uniformly (Appendix~\ref{app:baselines}): each model
occupies a \emph{level} on each of four tiers---architecture, rate, channels/placement, and label
vocabulary---and we never modify a released checkpoint's input contract, because doing so
discards the model being tested. CrossHAR and LiMU-BERT release no checkpoint compatible with our
corpus, so we pretrain them on \emph{our} data with their own published objectives, making their
data identical to ours; harnet, UniMTS and NormWear are used as released, which is faithful but
means their pretraining corpora differ from ours and from each other. This distinction is
material---harnet is pretrained on roughly $700$k person-days against our $572$
hours---and we state it wherever a comparison could be read as data-matched. Closed-vocabulary
models are bridged to arbitrary label strings with ConSE~\cite{conse}, so that ``can predict an
unseen label'' is available to every model and is not something HALO is credited for.

\subsubsection{Protocol}
\label{sec:eval_protocol}

\emph{Scoring.} Following the zero-shot evaluation literature~\cite{zslgbu}, every held-out
dataset is scored against its \emph{own} label strings, never a shared training-wide ontology, so
a prediction counts only if it retrieves that dataset's exact activity name. Macro-F1 is primary,
with accuracy alongside, because held-out label distributions are imbalanced.

\emph{Splits.} All few-shot fits and memory writes use subjects disjoint from those evaluated,
rotating over grouped folds with a confidence interval where a cohort is small enough that one
split is unstable. TNDA-HAR is distributed without per-window subject identifiers, so a
subject-disjoint boundary cannot be constructed for it; we report it separately and exclude it
from the subject-disjoint mean rather than implying a guarantee we cannot honour.

\emph{Label leakage.} The inference-time label-text ensemble is purged of synonyms drawn from the
evaluation datasets, so no evaluation vocabulary enters through paraphrase, and retrieval
hyper-parameters are selected on held-out \emph{training} configurations.

\emph{Parity control.} To separate the architecture from its native-rate and rich-text inputs,
HALO is also evaluated under baseline-equivalent inputs---$20$\,Hz resampling, neutral sensor
description---against the best member of each baseline group (Appendix~\ref{app:parity}).

\subsection{Q1: The Label-Efficiency Curve}
\label{sec:label_efficiency}

This is the headline experiment (Table~\ref{tab:label_efficiency}). At each budget $k$ of
labelled target windows per class we compare three mechanisms on identical data: HALO's memory, given the exemplars as an append with
\emph{no gradient update}; a probe on HALO's frozen features; and the same probe on frozen harnet
features. We additionally report \emph{retrieval over harnet features}---the same mechanism on a
different encoder---because without it a win could be credited to the mechanism when it belongs
to the representation, or the reverse.

The budget is reported in two units, because neither alone is honest: per-class $k$ is what a
deployment controls, while the fraction of the target dataset makes the numbers comparable to the
$1\%$/$5\%$/$10\%$ convention of the transfer literature. On our held-out sets the two coincide
closely---$k{=}10$ per class is $0.7$--$4.8\%$ of a dataset---and the grid stops at $k{=}100$
because beyond it the budget exceeds what the smaller sets contain.

\begin{table}[t]
  \centering
  \caption{Macro-F1 versus label budget, averaged over the held-out datasets, subject-disjoint
  throughout. $k{=}0$ is the zero-shot point and $k{=}\text{full}$ the supervised one; they are
  ends of one curve, not separate experiments, and every row uses \emph{one} frozen encoder
  across all $k$. The second header row converts $k$ into the mean fraction of a target dataset
  it represents. $^\dagger$\,The mechanism-versus-representation control: if retrieval helps only
  on HALO features the contribution is the encoder; if it helps equally on harnet's, it is the
  memory. We pre-register that we report this control regardless of which way it falls.}
  \label{tab:label_efficiency}
  \scriptsize
  \setlength{\tabcolsep}{1.5pt}
  \begin{tabular}{@{}l cccccc@{}}
    \toprule
    & \multicolumn{6}{c}{labelled target windows per class ($k$)} \\
    \cmidrule(lr){2-7}
    Mechanism & 0 & 5 & 10 & 25 & 100 & full \\
    \emph{(mean \% of target)} & \emph{0} & \emph{1.3} & \emph{2.7} & \emph{6.6} & \emph{27} & \emph{100} \\
    \midrule
    HALO memory (no gradient step)   & \PH & \PH & \PH & \PH & \PH & \PH \\
    HALO frozen $+$ probe            & \PH & \PH & \PH & \PH & \PH & \PH \\
    harnet frozen $+$ probe          & \PH & \PH & \PH & \PH & \PH & \PH \\
    harnet frozen $+$ memory\rlap{$^\dagger$} & \PH & \PH & \PH & \PH & \PH & \PH \\
    \midrule
    UniMTS (native text)             & \PH & \PH & \PH & \PH & \PH & \PH \\
    CrossHAR $+$ ConSE               & \PH & \PH & \PH & \PH & \PH & \PH \\
    LiMU-BERT $+$ ConSE              & \PH & \PH & \PH & \PH & \PH & \PH \\
    NormWear (native text)           & \PH & \PH & \PH & \PH & \PH & \PH \\
    \bottomrule
  \end{tabular}
\end{table}

\textbf{Pre-registered hypothesis and kill criterion.} We expect retrieval to win at small $k$,
where a nearest-neighbour decision beats a linear fit on a handful of points, and the probe to
overtake somewhere between $k{=}10$ and $k{=}25$---roughly the $3$--$7\%$ band. If that holds,
the regime HALO owns is $k \le 10$, which is the enrollment setting a user can perform in under a
minute. If the probe dominates at \emph{every} $k$ including the smallest, retrieval has no
regime of its own; we will report that outcome plainly. Appendix~\ref{app:perdataset} reports the
per-dataset breakdown, including InclusiveHAR stratified by ability, because a model that works
only for able-bodied users is not a foundation for rehabilitation applications.

\subsection{Q2: Does the Inductive Bias Do What It Claims?}
\label{sec:bias_experiments}

Section~\ref{sec:no_grammar} made two falsifiable claims about the tokenizer. Both are tested
directly rather than inferred from end-task accuracy.

\textbf{Rate-invariance is measured, not assumed.} A closed-form check confirms a pure tone
deposits identical band energies at $20$, $50$, and $100$\,Hz. The end-task version is
Table~\ref{tab:tokenizer_ablation}: the same encoder on the same data but reading a common
$20$\,Hz resampling of every stream. If the faster streams carry discriminative content above
$10$\,Hz, resampling must destroy it, and the gap measures how much.

\textbf{Does anchoring actually make one cross-device memory well-posed?} This experiment decides
whether HALO is an assembly of known parts or an enabling result. A flat memory shared across
devices is only meaningful if a query encoded from one configuration retrieves useful neighbours
stored from a \emph{different} one, which is what Table~\ref{tab:crossconfig} measures. The
measurement is unusually clean because parts of our corpus contain \emph{verified simultaneous}
recordings---the same person, activity, and instant, captured by two or more devices at different
placements---so querying one stream against a memory built only from its simultaneous partner
isolates configuration from subject, session, and activity entirely. We report the same matrix
for a resample-to-common-rate variant of our own encoder and for the frozen baselines, so the
comparison is of mechanisms rather than corpora.

\textbf{Is the naive assembly enough?} Since HALO's components have individually appeared before,
the claim that they must be assembled \emph{this particular way} is empirical.
Table~\ref{tab:tokenizer_ablation} includes the obvious alternative as a row: resample every
stream to a common rate, attach free text per channel, same retrieval on top. If that arm matches
HALO the architecture is a combination and we will say so; mutual entailment predicts it will
not, and predicts a reason---per-channel conditioning showing up as elevated dataset
identifiability, resampling as a rate gap.

\begin{table}[t]
  \centering
  \caption{\textbf{Cross-configuration retrieval: the enabling claim.} Each row queries with one
  acquisition configuration against a memory containing \emph{only} a different one. Because
  these pairs are verified simultaneous recordings, the only thing that varies is the
  configuration, so off-diagonal degradation is attributable to it alone. \emph{Retention} is
  cross divided by same: features that denote the same physical quantity across devices should
  retain most of their quality; one that resamples or conditions on the rate should not.}
  \label{tab:crossconfig}
  \scriptsize
  \setlength{\tabcolsep}{4pt}
  \begin{tabular}{@{}l cc cc@{}}
    \toprule
    & \multicolumn{2}{c}{HALO} & \multicolumn{2}{c}{resample-to-20\,Hz} \\
    \cmidrule(lr){2-3}\cmidrule(lr){4-5}
    Query $\rightarrow$ memory & same & \textbf{cross} & same & \textbf{cross} \\
    \midrule
    wrist $\rightarrow$ pocket (SP-SW-HAR)   & \PH & \PH & \PH & \PH \\
    wrist $\rightarrow$ lower back (NFI)     & \PH & \PH & \PH & \PH \\
    pocket $\rightarrow$ watch (WISDM)       & \PH & \PH & \PH & \PH \\
    wrist $\rightarrow$ ear/glasses (XRF~V2) & \PH & \PH & \PH & \PH \\
    100\,Hz $\rightarrow$ 20\,Hz memory      & \PH & \PH & \PH & \PH \\
    \midrule
    \textbf{cross/same retention}            & \multicolumn{2}{c}{\PH} & \multicolumn{2}{c}{\PH} \\
    \bottomrule
  \end{tabular}
\end{table}

\textbf{Conditioning must not be a corpus fingerprint.} The right-hand columns of
Table~\ref{tab:tokenizer_ablation} report a dataset-identity probe: a linear classifier fitted on
frozen embeddings of held-out subjects to predict which \emph{source dataset} a window came from.
A representation that has absorbed the acquisition description as a corpus identifier makes this
easy; one whose conditioning is genuinely compositional does not. It is the only direct evidence
we have for the ``too little bias'' failure mode, so we report it whether or not it favours us.

\begin{table}[t]
  \centering
  \caption{Tokenizer inductive-bias ablation (macro-F1 over the held-out datasets), each
  component removed independently, with the dataset-identity probe alongside. The
  \emph{naive assembly} row is not an ablation but a \emph{rival}: the same retrieval mechanism
  over an encoder built the conventional way. It is the arm that tests whether HALO is a
  combination of known parts. $^\ddagger$\,\emph{Rate gap} is performance on the $100$\,Hz
  held-out cell minus the same data downsampled to $20$\,Hz---what resampling costs.
  \emph{ds-ID} is dataset-identity accuracy; lower is better.}
  \label{tab:tokenizer_ablation}
  \scriptsize
  \setlength{\tabcolsep}{1.5pt}
  \begin{tabular}{@{}l cc cc@{}}
    \toprule
    & \multicolumn{2}{c}{transfer} & \multicolumn{2}{c}{probe} \\
    \cmidrule(lr){2-3}\cmidrule(lr){4-5}
    Tokenizer configuration & $k{=}0$ & $k{=}10$ & rate gap\rlap{$^\ddagger$} & ds-ID $\downarrow$ \\
    \midrule
    Physical-Hz, sensor-grouped, factored text & \PH & \PH & \PH & \PH \\
    \quad$-$ filterbank (resample to 20\,Hz)   & \PH & \PH & \PH & \PH \\
    \quad$-$ sensor grouping (per-chan.\ text) & \PH & \PH & \PH & \PH \\
    \quad$-$ language conditioning (locked)    & \PH & \PH & \PH & \PH \\
    \quad$-$ observability / resolution masks  & \PH & \PH & \PH & \PH \\
    \quad$-$ signed DC (gravity) feature       & \PH & \PH & \PH & \PH \\
    \midrule
    \emph{naive}: resample $+$ per-chan.\ text & \PH & \PH & \PH & \PH \\
    \midrule
    harnet~\cite{sslwearables} / UniMTS~\cite{unimts} & \PH & \PH & --- & \PH \\
    chance                                     & --- & --- & --- & \PH \\
    \bottomrule
  \end{tabular}
\end{table}

\subsection{Q3: Which Phase-A Objectives Carry the Transfer?}
\label{sec:objective_ablation}

Each Phase-A term is removed independently with all others held fixed
(Appendix~\ref{app:objectives}). Two controls matter more than the ablation itself. The
\emph{untrained control}: the decoder is identity-at-initialisation by construction, so every
trained arm is reported against the same architecture untrained, and a component that does not
beat its own initialisation is reported as such. The \emph{label-supervised control}: the
original supervised recipe remains selectable, so the claim that going label-free helps is
measured, not asserted.

\subsection{Q4: On-Device Cost and a Real Deployment}
\label{sec:deployment_experiments}
\label{sec:case_study}

We implement the full pipeline as an iOS application and will report encoder latency and peak
memory under Core~ML, retrieval latency against bank size, and a case study in which a user
enrolls activities by demonstration, changes body placement mid-session, and introduces names
never enrolled---the deployment form of the label-efficiency argument.

\PENDING{On-device latency, peak memory, retrieval cost versus bank size, and the case-study
figure are added here once the encoder is compiled and profiled.}
